\documentclass[11pt]{article}
\usepackage[margin=1in]{geometry}
\usepackage{amsmath,amssymb,booktabs}
\usepackage[hidelinks]{hyperref}
\title{Quantum control, formation, covariant response, and vacuum normalization}
\author{Zimmerman programme: CPB-QFG-2026-10-16}
\date{16 October 2026}
\begin{document}
\maketitle
\section{Results and their limits}
Four requested gaps are tested on the actual mediator candidate. Its
original order-one couplings do not protect the small attractive quartic
or the light host mass under one-loop running. A separately specified
weak-coupling scaling suppresses those sensitivities while preserving
classical host shapes and instability rates, at the cost of larger charge
and energy inventories. Dilute charged coherent backgrounds have a linear
clumping instability; this does not derive a cosmological charged initial
state or nonlinear droplet production.

The minimal Lorentz-invariant scalar realization of the cubic static law
has a radial characteristic speed approaching $\sqrt2$ times the metric
light speed. Thus the static force calculation alone does not supply a
completion with a common metric cone. Independent medium dynamics may
change that result but have not yet been derived. Finally, the present
particle and gravity equations leave independent response coefficients and
vacuum flux/history data. They do not predict a unique dark-energy density
or acceleration normalization. None of these gaps is relabeled as closed.

\section{Full scalar one-loop divergence and matching sensitivity}
Write the mediator potential in general form
\begin{equation}
 V=V_0+t\chi+m_h^2 I+L I^2+a\chi I+
 \tfrac12 B\chi^2+\tfrac\mu3\chi^3+\tfrac l4\chi^4,
 \qquad I=|\Psi|^2,\quad B=M^2.
\end{equation}
The previous tree family has $V_0=t=0$, $m_h=1$,
$L=2-1/M$, $a=2M$, $\mu=-3M/8$, $l=1$.
The independent vacuum and tadpole counterterms must nevertheless be kept.
For $\Psi=(\phi_1+i\phi_2)/\sqrt2$ and
$\phi_1=\sqrt{2I},\phi_2=0$, the canonical real-scalar Hessian is
\begin{equation}
 H=\begin{pmatrix}
 m_h^2+6LI+a\chi&0&a\sqrt{2I}\\
 0&m_h^2+2LI+a\chi&0\\
 a\sqrt{2I}&0&B+2\mu\chi+3l\chi^2
 \end{pmatrix}.
\end{equation}
The real-scalar MS one-loop convention is
\begin{equation}
 V_1=\frac{1}{64\pi^2}\sum_j h_j^2
       [\ln(h_j/\mu_R^2)-3/2],\qquad h_j=\operatorname{eig}(H).
\end{equation}
Its normalization is checked against Ref.~\cite{qfg:Martin}; the following
polynomial projections are derived here. This expression for a stable
background is not a finite-density fluctuation determinant at nonzero
chemical potential. Negative $h_j$ require the usual instability treatment;
no real global quantum potential is inferred from them.

At one loop the ultraviolet field-strength poles vanish for these scalar
interactions. Scale independence gives
$\sum_i\beta_i\partial_i V=\operatorname{tr}H^2/(32\pi^2)$.
The exact polynomial is
\begin{align}
 \operatorname{tr}H^2={}&40L^2I^2+(16Lm_h^2+4a^2)I+16LaI\chi
 +9l^2\chi^4+12l\mu\chi^3\nonumber\\
 &+(2a^2+6Bl+4\mu^2)\chi^2
 +(4m_h^2a+4B\mu)\chi+2m_h^4+B^2.
 \label{qfg:trace}
\end{align}
Let $\ell_1=16\pi^2$. Its coefficients yield
\begin{align}
 \ell_1\beta_L&=20L^2,&\ell_1\beta_a&=8La,
 &\ell_1\beta_{m_h^2}&=8Lm_h^2+2a^2,\\
 \ell_1\beta_B&=2a^2+6Bl+4\mu^2,
 &\ell_1\beta_\mu&=18l\mu,&\ell_1\beta_l&=18l^2,\\
 \ell_1\beta_t&=2m_h^2a+2B\mu,
 &\beta_{V_0}&=(2m_h^4+B^2)/(32\pi^2).
\end{align}
These are full-theory MS beta functions at this potential's zero quartic
mixed-operator point, not finite threshold matching or running in a
heavy-field-decoupled low-energy EFT.

For the attractive tree coefficient $\lambda_{\rm net}=L-a^2/(2B)$,
\begin{equation}
 \ell_1\beta_{\lambda_{\rm net}}
 =20L^2-\frac{8La^2}{B}
 +\frac{a^2}{2B^2}(2a^2+6Bl+4\mu^2).
\end{equation}
On the selected family this becomes
$[20L^2-32L+233/8]/\ell_1$.
At $M=100$, $\lambda_{\rm net}=-0.01$ but
$\beta_{\lambda_{\rm net}}=0.2827304$;
$\beta_{m_h^2}/m_h^2=506.7067$. The small matching relation is not protected
by the tree vacuum proof. A linearized scale interval
$|\delta\ln\mu_R|\simeq0.03537$ already corresponds to a shift equal to
the chosen quartic magnitude. This diagnostic is not the claim that a
physical observable varies with an arbitrary renormalization scale:
consistent loop matching restores scale invariance, and finite
renormalization can impose measured inputs at a chosen scale.

Tadpoles are not omitted. With $t=0$ initially, following the stationary
vacuum gives $dv/d\ln\mu_R=-\beta_t/B$. At that point the corresponding
curvature-mass and local tree-quartic projections are
\begin{align}
 \beta_{m_h^2}^{\rm follow}&=\beta_{m_h^2}-a\beta_t/B,\\
 \beta_{\lambda_{\rm net}}^{\rm follow}
 &=\beta_{\lambda_{\rm net}}-a^2\mu\beta_t/B^3.
\end{align}
On the family the second correction is
$(6/M^2-9/8)/\ell_1$. It does not remove the order-one-coupling sensitivity.
These local projections are still not pole-mass or scattering threshold
matching conditions.

\section{A controlled classical-limit route}
Introduce a new freely selected parameter $0<\epsilon\ll1$ and replace
\begin{equation}
 L\mapsto\epsilon L,\quad a\mapsto\sqrt\epsilon a,\quad
 \mu\mapsto\sqrt\epsilon\mu,\quad l\mapsto\epsilon l,
 \quad m_h^2,B\ \hbox{fixed}.
\end{equation}
For $\Psi=\psi/\sqrt\epsilon$ and $\chi=X/\sqrt\epsilon$, the exact
scalar action is $S_\epsilon=S_1[\psi,X]/\epsilon$.
Its classical field equations, background Hessian eigenvalues and
formation dispersion relations in the rescaled fields are identical.
The stable-vacuum proof and coexistence phase persist under this invertible
rescaling. Charge, energy and interface tension for a fixed shape grow as
$1/\epsilon$; the radius does not grow unless the specified inventory does.

The matching coefficient scales as $\epsilon\lambda_{\rm net}$,
its one-loop beta function as $\epsilon^2\beta_{\lambda_{\rm net}}$,
and the mass beta as $\epsilon\beta_{m_h^2}$. Consequently a weakly coupled
version can control both diagnostics without changing the classical
host mechanism. At $M=100$:
\begin{center}
\begin{tabular}{rrrr}
\toprule
$\epsilon$ & $|\beta_{\lambda_{\rm net}}/\lambda_{\rm net}|$
 & $|\beta_{m_h^2}/m_h^2|$ & inventory multiplier\\
\midrule
1 & 28.273 & 506.707 & 1\\
$10^{-3}$ & 0.02827 & 0.50671 & $10^3$\\
$10^{-6}$ & $2.8273\times10^{-5}$ & $5.0671\times10^{-4}$ & $10^6$\\
\bottomrule
\end{tabular}
\end{center}
The last row passes a specified one-percent, unit-log sensitivity check;
the first two do not. This is a quantum-control route, not a calculation
of all finite thresholds or a quantum vacuum theorem. Large classical
fields require control of any additional operators in a higher theory.
Portal interactions, response couplings, gravitational backreaction and
cosmological inventories must also be matched anew. The parameter
$\epsilon$ and the increased stock have not been predicted.

\section{Formation: an actual instability, with specified initial state}
At a homogeneous charged background, solve the actual mediator equation
$U'(\chi)+aI=0$ and
$\omega^2=m_h^2+2LI+a\chi>0$. Write
$\Psi=e^{i\omega t}[\sqrt I+(\eta+i\pi)/\sqrt2]$ and perturb $\chi$
by $\xi$. Define
\begin{equation}
 A=4LI,\qquad B_h=U''(\chi),\qquad h=a\sqrt{2I},\qquad
 A_{\rm eff}=A-h^2/B_h.
\end{equation}
For wave number $k$, the exact linear scalar system has potential matrix
\begin{equation}
 K(k)=\begin{pmatrix}k^2+A&0&h\\0&k^2&0\\h&0&k^2+B_h\end{pmatrix}
\end{equation}
and antisymmetric rotation coupling $G_{12}=-G_{21}=-2\omega$ in
$\ddot q+G\dot q+Kq=0$. For $z=\Omega^2$, its dispersion equation is
\begin{equation}
 (k^2+B_h-z)[(k^2+A-z)(k^2-z)-4\omega^2 z]
 -h^2(k^2-z)=0.
 \label{qfg:dispersion}
\end{equation}
The long-wavelength mode has
$c_s^2=A_{\rm eff}/(A_{\rm eff}+4\omega^2)$.
If $A_{\rm eff}<0$ and the denominator is positive, there is a growing
long-wavelength mode. Its unstable band ends at
\begin{equation}
 k_c^2=\frac{2(h^2-AB_h)}
 {\sqrt{(A-B_h)^2+4h^2}+A+B_h}>0.
\end{equation}
A positive $A_{\rm eff}$ instead makes the quadratic energy nonnegative
for all $k$, giving linear stability in this conservative scalar sector.

The real-time six-dimensional generator was diagonalized at 201 wave
numbers $0\le k\le2$, for $M=10,100$ and $I/I_*=0.1,0.5,1,2$.
Its eigenvalues agree with the independently derived stable/unstable-band
classification across the sampled points. At $M=100$:
\begin{center}
\begin{tabular}{rrrrr}
\toprule
$I/I_*$ & $c_s^2$ & $k_c$ & fastest sampled $k$ & growth rate $\gamma$\\
\midrule
0.1 & $-0.04815$ & 0.40795 & 0.29 & 0.04369\\
0.5 & $-0.09924$ & 0.49502 & 0.34 & 0.07436\\
1 & 0.24898 & 0 & --- & no sampled growth\\
2 & 0.50357 & 0 & --- & no sampled growth\\
\bottomrule
\end{tabular}
\end{center}
Units are $m_h$ for wave number/growth and $m_h^{-1}$ for time.
The weak-coupling scaling preserves these rates at the rescaled density.
This is a concrete possible fragmentation starting point, conditional on
preparing a charged coherent state in that range.

The empty vacuum has positive scalar masses and no such instability. A
neutral thermal gas is not this coherent charged state. Its production,
charge separation and occupation cannot be inferred from the table.
Nonlinear evolution, charge/energy conservation and radiation must be
executed before claiming actual droplets or abundance. An isolated system
can radiate into its own scalar waves; an external heat bath is not
logically necessary. The calculation does not establish which available
radiation channels can remove the binding energy at a sufficient rate.
The raw outputs record the initial energy per charge and the bulk binding
budget; finite surfaces and emitted charge must be counted separately.

\section{A covariant coupling test: the minimal scalar cone}
A conformal baryon metric $\widetilde g_{\mu\nu}=e^{2q_b\varphi}g_{\mu\nu}$
reproduces the weak static source energy $q_b\rho_b\varphi$. A Lorentz
auxiliary vector can reproduce the static square through
\begin{equation}
 \mathcal L_{\rm resp}=-\tfrac C2(\partial\varphi)^2
 +P^\mu\partial_\mu\varphi-\frac{P^\mu P_\mu}{2C}-W(P^2).
\end{equation}
On its spacelike static branch choose $W=\beta(P^2)^{3/2}/3$.
This defines a local branch, not a prescribed timelike or vacuum extension.
Eliminating $P$ yields a single-gradient scalar Lagrangian $-F(X)$,
$X=(\partial\varphi)^2/2>0$. If $y=|\nabla\varphi|$, the exact
constitutive equations are
\begin{equation}
 y=p/C+\beta p^2,\qquad F_y=D=C\beta p^2.
\end{equation}
The principal tensor is
$F_Xg^{\mu\nu}+F_{XX}\partial^\mu\varphi\partial^\nu\varphi$.
On the static radial branch it gives
\begin{equation}
 c_r^2=\frac{y\,dD/dy}{D}
 =\frac{2(1+C\beta p)}{1+2C\beta p},\qquad c_\perp^2=1.
 \label{qfg:cone}
\end{equation}
Thus the cubic-dominated limit $C\beta p\ll1$ has $c_r^2\to2$.
The temporal kinetic coefficient is positive for $p>0$ but tends to zero
as the critical branch approaches $p=0$; a uniform perturbative vacuum
limit is not supplied. Conformal coupling leaves the baryonic null cone
unchanged, so this is faster than that metric cone. It fails a requirement
that every mode propagate within the same cone. Superluminal scalar
characteristics alone are not a complete proof of a causal contradiction
or an empirical exclusion of all alternatives.

For the quadratic exterior $W=K_0P^2/2$, the vector can be eliminated
covariantly without a material frame and gives a positive scalar kinetic
coefficient $C_{\rm eff}=C^2K_0/(1+CK_0)$. The exterior therefore need
not rely on the undefined phase current outside a droplet. Its enhanced
Newtonian force remains the earlier $1+1/(CK_0)$ factor.
Independent medium/frame dynamics may alter the interior characteristic
relation; their action and radiative matching still have to be derived.
Nothing in the scalar mediator potential fixes $\beta$, $C$, $q_b$ or the
critical response switching functions. Relativistic lensing and the full
metric field equations remain required as well.

\section{Dark energy: identifiable inputs and a decisive degeneracy}
The local one-loop vacuum divergence includes
$(2m_h^4+M^4)/(32\pi^2)$ in $\beta_{V_0}$, with an independent finite
vacuum counterterm. It is not a predicted cosmological constant. Even the
weak-coupling limit does not remove the free massive vacuum determinant;
its relative smallness compared with a large classical condensate is not
an explanation of the observed vacuum scale.

The previous auxiliary construction instead has
\begin{equation}
 B_{\rm eff}=Q^2/2+\Delta_F+\langle\tau\rangle/4.
\end{equation}
Vacuum-form matter constants can project out of its metric equation under
its stated hypotheses, but their $Q$ derivatives remain in the conjugate
flux equation. That algebra does not select $Q$ or the residual
flux/history term. These are global data, not a result of the host's
local coexistence condition.

For $W=\beta p^3/3$, the annular force coefficient is
\begin{equation}
 a_0=\frac{q_b^3}{4\pi C^3\beta G_{\rm lab}}.
\end{equation}
Holding the particle potential and all $P=0$ phase data fixed while
changing $\beta$ from 1 to 4 changes $a_0$ by a factor of four.
Holding $Q=1$ fixed while changing
$\Delta_F+\langle\tau\rangle/4$ from 0 to 1 changes $B_{\rm eff}$
from $1/2$ to $3/2$. Both variations are allowed by the current parameter
and boundary-data specification. Therefore no unique $\kappa$ or physical
dark-energy normalization follows from it. These counterexamples apply
to this specification, not every possible future completion.

A further mechanism must fix the response coefficients and flux/history
data independently of the measured target. Otherwise the honest result
is a calibrated effective theory. Selecting those data from the observed
$a_0$ and dark-energy density is not a derivation. Flux quantization or a
vacuum-selection proposal would also need its own charge scale and actual
global solution; naming it supplies neither condition.

\section{Evidence and next implication}
\texttt{four\_gap\_audit\_run\_003} has 3219 exact/finite controls,
including 1608 generator spectra and independent band checks. Run 001
completed the core; run 002 added tadpole-following and energy budgets;
run 003 added independent spectral-band controls. No physical thresholds
were loosened. \texttt{weak\_coupling\_control\_run\_001} has five exact
scaling checks and three finite sensitivity comparisons with passing and
failing scientific flags retained. Both current manifests passed freshness
validation. These are self-reviewed results, not independent certificates.

The next useful joined calculation is nonlinear charged-state fragmentation
in the weakly coupled mediator model with explicit outgoing scalar energy
and charge, followed by a physically specified response/metric completion
that meets its declared characteristic-cone and lensing requirements.
Quantum finite matching and independently selected global vacuum data
remain necessary. The full cold-energy/dark-energy theory is still open.

\begin{thebibliography}{9}
\bibitem{qfg:Martin} S. P. Martin,
\emph{Two-loop effective potential for a general renormalizable theory and
softly broken supersymmetry}, arXiv:hep-ph/0111209v2,
\url{https://arxiv.org/pdf/hep-ph/0111209v2}, equations (1.1), (1.3), (3.2).
The previously authenticated real-scalar one-loop convention is reused;
new beta projections, dispersion and response characteristics are derived
above. No external finite-density or full-model stability theorem is imported.
\end{thebibliography}
\end{document}
